\chapter{Diretrizes operacionais}

\section{Estrutura da Coordenação de Comunicação}

\subsection{Atribuições}

À Coordenação de Comunicação do CRP~SP cabe a gestão de conteúdo do \emph{site}, dos perfis sociais, dos produtos fixos e das publicações diárias da instituição.

Também é responsabilidade da Coordenação de Comunicação manter boas relações com a imprensa, promovendo contato com jornalistas e veículos, acompanhando a realização de entrevistas e a publicação de matérias, elaborando notas de elucidação e preparando material de divulgação na forma de \emph{press release} sempre que necessário.

\subsection{Responsabilidades sobre gestão e administração}

\begin{itemize}
	\item
	Passagens, estadias e outros custos de escopo estadual;
	\item
	administração financeira do setor (orçamentos, termos de licitação,
	atestes);
	\item
	planejamento estratégico da área.
\end{itemize}

\subsection{Recursos}

\section{Processos e produtos comunicacionais}

\subsection{Canais oficiais}

\subsubsection{\emph{Site} institucional}

\subsubsection{Redes sociais}

O CRP~SP mantém perfis ativos nas seguintes redes sociais:

\begin{itemize}
	\item
	Instagram;
	\item
	Facebook;
	\item
	LinkedIn;
	\item
	X/Twitter;
	\item
	YouTube.
\end{itemize}

Os perfis sociais do CRP~SP têm como objetivo formar e informar a categoria sobre temas associados à Psicologia, à atuação profissonal e ao funcionamento do Sistema Conselhos. Por meio deles, buscamos ampliar nossos pontos de contato e comunicação com as psicólogas e psicólogos de todo o estado de São Paulo, além de participar dos debates sobre a
Psicologia como ciência e profissão.

Em observância aos princípios da transparência e do compromisso ético, evidenciamos as seguintes diretrizes:

\begin{enumerate}
	\def\labelenumi{\arabic{enumi}.}
	\item
	\textbf{perfis e páginas oficiais}: só são permitidos perfis e páginas
	associados ao CRP~SP que sejam criados e mantidos pela instituição. A
	autarquia se reserva o direito de não endossar perfis e páginas que se
	autointitulem associados ao Conselho e que não sejam reconhecidos como
	oficiais, tomando as devidas medidas quando necessário;
	\item
	\textbf{gestão das contas}: as contas oficiais do CRP~SP serão
	gerenciadas por uma equipe certificada e autorizada, composta por
	profissionais da comunicação, trabalhadores da autarquia e psicólogas
	e psicólogos conselheiros;
	\item
	\textbf{política de interação e moderação}: para garantir a
	veracidade, exatidão, confiabilidade, qualidade e legalidade das
	informações, as interações nos canais oficiais estarão sujeitas a
	moderação. Mensagens e perfis poderão ser excluídos se:

	\begin{enumerate}
		\def\labelenumii{\arabic{enumii}.}
		\item
		usarem linguagem inapropriada, obscena, caluniosa, grosseira,     abusiva, difamatória ou ofensiva;
		\item
		fizerem apologia a práticas ilícitas;
		\item
		incitarem ódio, violência, racismo ou qualquer forma de discriminação;
		\item
		contiverem ameaças, assédio, injúria, calúnia, difamação ou qualquer outro ilícito penal;
		\item
		divulgarem conteúdos na forma de \emph{spam} ou~``correntes'';
		\item
		caracterizarem intuito comercial ou publicitário;
		\item
		forem ininteligíveis ou fora de contexto;
		\item
		contiverem propagandas político-partidárias;
		\item
		contiverem \emph{links} suspeitos ou representarem ameaça à segurança da informação;
		\item
		usarem informações e imagens de pessoas e instituições de modo indevido;
		\item
		contiverem dados pessoais da autora, do autor ou de terceiros;
		\item
		violarem direitos de imagem e de propriedade intelectual;
		\item
		forem fraudulentos (perfis falsos e/ou \emph{bots}) ou promoverem conteúdo inverídico (\emph{fake news}).\\
	\end{enumerate}
	\item
	\textbf{restrições de divulgação}: o CRP~SP não segue e não marca pessoas, veículos de comunicação e empresas privadas, apenas entidades parceiras e órgãos do poder público. Não divulgamos partidos políticos e não realizamos atendimento administrativo e de orientação nas redes sociais. As demandas fora dos produtos fixos devem ser levadas a
	comissões estaduais.
\end{enumerate}

A não observância dessas diretrizes pode acarretar ações imediatas do CRP~SP, sem necessidade de justificativa, consulta ou aviso prévio. Mensagens inadequadas poderão ser encaminhadas às autoridades responsáveis. Nosso compromisso é com a ética, a transparência e a qualidade das nossas comunicações.

\subsection{Produtos fixos}

São produtos de comunicação do CRP~SP:

\subsubsection{Jornal Psi}

Publicado em quatro edições anuais de maneira \emph{on-line}, com envio de versão impressa a psicólogas, psicólogos e entidades com inscrição ativa no CRP~SP. Pautas e validação final do jornal são compartilhadas com o Plenário.

\subsubsection{Psicologia em Ação}

Divulgação de eventos a pedido, utilizando-se do \emph{site} (Agenda CRPSP), de \emph{stories} (perfil oficial no Instagram), postagens (demais redes sociais) e \emph{cards} a serem compartilhados pelos territórios. A cobertura dos eventos é publicada no CRP~SP Acontece.

As publicações de \emph{Psicologia em Ação} possuem identidade visual própria, que pode receber elementos temáticos.

\subsubsection{CRP~SP Acontece}

Veiculado semanal ou quinzenalmente, noticia a categoria e a sociedade sobre as ações, eventos, representações, mobilizações, novidades etc. do dia a dia do CRP~SP, em formato de galeria de imagem (até 10 fotos por publicação) acompanhada por breves relatos.

Subsedes , comissões e demais instâncias devem enviar seus reportes (breves: um parágrafo contando o que foi a ação, quem esteve presente etc.), acompanhados de foto para Comunicação.

\subsubsection{CRP~SP Informa}

Ferramenta de \emph{e-mail marketing} destinada a dar publicidade às ações realizadas nos territórios, por meio de envio de mensagens institucionais a psicólogas, psicólogos e entidades afins. É uma importante ferramenta de comunicação direcionada, com média estadual de abertura dos \emph{e-mails} de 20\%.

A cada território é facultado o envio de até quatro mensagens por mês.

\subsubsection{CRP~SP Debate (\emph{lives})}

Canal que possibilita ampliar discussões para mais públicos e trazer convidadas de diferentes espaços. As \emph{lives} são realizadas por comissões e, geralmente, estão inseridas em uma campanha. Contam sempre com a tradução para Libras e têm o \emph{chat} mediado pela Coordenação de Comunicação. São operadas no StreamYard (transmissão pelo YouTube e Facebook), podendo apresentar até 7 participantes.

\emph{Lives} devem ser solicitadas com mínimo de 30 dias de antecedência, por meio de formulário de eventos, que deve ser preenchido e enviado para as pessoas responsáveis por eventos na Coordenação de Comunicação. Após o recebimento, será criado grupo de imMail para realização de testes prévios, elaboração de roteiro e alinhamento necessário.

\subsubsection{Calendário de Direitos Humanos}

Artigos e postagens publicados em datas simbólicas para a defesa dos Direitos Humanos. Textos são organizados pelas comissões permanentes e encaminhados à ComCom com até três dias úteis de antecedência para publicação, e devem ser elaborados com linguagem inclusiva (\#paratodosverem ), fornecendo informações sobre o significado da data em destaque, abordando aspectos relevantes do exercício profissional da Psicologia e fazendo referência à(s) normativa(s) do Sistema Conselhos e/ou do Código de Ética alinhada(s) ao assunto da data. A escrita deve ser objetiva, simples e compatível com as redes sociais (máximo de dois mil caracteres e quatro parágrafos). O conteúdo é padronizado e institucional.

\subsubsection{CRP~SP na Mídia}

Divulga as participações do CRP~SP na imprensa e em outros canais externos.

Solicitações devem ser encaminhadas para o \emph{e-mail} \nolinkurl{ascom@crpsp.org.br}. Os pedidos são levados para representações e
plenário, para que se decida sobre o encaminhamento.

\subsection{Publicações diárias}

A Coordenação de Comunicação do CRP~SP é responsável pela publicação a qualquer tempo, de acordo com a demanda, de:

\begin{itemize}
	\item
	material do Sistema Conselhos de Psicologia;
	\item
	orientações do Centro de Referências Técnicas em Psicologia e Políticas Públicas (Crepop);
	\item
	eventos e outros assuntos de interesse produzidos por entidades parceiras;
	\item
	notas de pesar;
	\item
	notas de posicionamento;
	\item
	informes administrativos e financeiros;
	\item
	orientações gerais.
\end{itemize}

\subsection{Campanhas}

São solicitadas por comissões permanentes e especiais e unidades administrativas. Contam com plano de comunicação elaborado em parceria com a ComCom trazendo o escopo completo da campanha (mote, objetivo, ações, grade de veiculação etc.).

\subsection{Eventos}

A Coordenação de Comunicação do CRP~SP é responsável pela gestão de eventos promovidos pela autarquia \textbf{em nível estadual}, como mostras, seminários, rodas de conversa, congressos, cinedebates, premiações e eventos internos.

A produção de eventos nos territórios é atribuição das respectivas subsedes, cabendo à Coordenação de Comunicação a orientação e a divulgação, além da disponibilização das seguintes ferramentas digitais:

\begin{itemize}
	\item
	inclusão do evento na Agenda CRP~SP;
	\item
	fornecimento de \emph{link} para inscrição;
	\item
	emissão de certificados.
\end{itemize}

\subsubsection{Atribuições das subsedes}

\subsubsection{Fluxo de eventos}

\paragraph{Requisição}

\begin{enumerate}
	\def\labelenumi{\arabic{enumi}.}
	\item
	Subsede preenche calendário de eventos na nuvem;
	\item
	envia o formulário de evento (roda de conversa \emph{on-line} ou   presencial), com informações completas, para as pessoas responsáveis por eventos na Coordenação de Comunicação;
	\item
	responsável abre inscrição na Agenda CRP~SP (no \emph{site}  institucional), verifica a necessidade de tradução para Libras, encerra as inscrições um dia antes do evento e envia a lista de inscritos para a subsede;
	\item
	Coordenação de Comunicação cria \emph{card} de divulgação, envia para validação da subsede por \emph{e-mail} e divulga nos canais oficiais;
	\item
	subsede envia listas de presença para que responsáveis por eventos liberem certificado de participação na Agenda CRP~SP.
\end{enumerate}

\paragraph{Produção}

\begin{itemize}
	\item
	Inscrições;
	\item
	certificados;
	\item
	divulgação;
	\item
	produção;
	\item
	plenária.
\end{itemize}

\subsubsection{Cobertura e transmissão}

\section{Gestão e proteção de
	dados}

\subsection{Audiovisual}

\subsubsection{Arquivamento}

Todo o material audiovisual produzido deve ser arquivado e indexado, para que a qualquer tempo seja possível a busca por meio dos seguintes dados:

\begin{enumerate}
	\def\labelenumi{\arabic{enumi}.}
	\item
	data de criação;
	\item
	descrição geral do assunto (nome de evento, campanha em que se insere  etc.);
	\item
	nomes das pessoas gravadas, filmadas ou fotografadas.
\end{enumerate}